\section{Memory}

A computer needs to store and recall, and the mesh has a natural memory built from the same parts.

A stored pattern is a stable arrangement of tones hung at an address on the growth tree. Because the tree gives every
cell a short name, a pattern can be placed at a known location and found again by navigating to that name. A still
pattern is a memory, a remembered thing read when a path crosses it, and a running pattern is a program, a precomposed
sequence that does not merely record an experience but replays it, driving whoever crosses it through the steps it
encodes.

The structure of the tree makes this memory unusually capable. The set of all cell names is at once a trie, general at
the root and specific at the leaves, a content-addressable store reachable by what a pattern is and not only by where it
sits, a balanced index of logarithmic depth, and an error-correcting code, because a pattern spread across many cells
survives the loss of some of them. An experiment confirms the error-correcting behavior: a spread-encoded bit survives a
bounded erasure while a localized blob is destroyed, so storing a pattern across the branching tree protects it. The
capacity is large because the hyperbolic interior is roomy, and the search is fast because the tree is shallow, the same
two facts that gave navigation and holography.

The one real limit is persistence. Whether a pattern written into the churning bulk lasts against the constant shuffling
of the rule, unless it is error-corrected, is open, the persistence problem. A pattern protected by spreading and
correction survives, as the experiment shows. Whether the bare rule sustains that correction on its own, long enough for
a stable self to persist, is not yet settled.
